% !PRINTED 195
\weberpage{195}
In this also lies the proof that, if $D$ is different from zero, the normal form can be produced by a linear transformation. For under this assumption $H$ decomposes into two linear factors distinct from one another, and if one chooses these, apart from constant factors, as the new variables $\xi$, $\eta$ of a linear transformation, then $A_0' = 0$, $A_2' = 0$; but it follows from the identical relations
\[
3\,a_3\,A_0 + a_1\,A_2 = a_2\,A_1,\quad a_2\,A_0 + 3\,a_0\,A_2 = a_1\,A_1,
\]
unless at the same time $A_1' = 0$, which is excluded by the non-vanishing $D$, that $a_1' = 0$, $a_2' = 0$, i.\,e. the transformed form of $f$ contains only the cubes of $\xi$ and $\eta$.

If one raises the first equation (10) to the cube, and eliminates $\xi^3$, $\eta^3$, $r^6$ by means of (11) and the second equation (10), one obtains among the covariants the following identical relation
\weq{12}{4\,H^3 + Q^2 + 27\,D\,f^2 = 0.}

To carry out the transformation into the normal form, i.\,e. actually to find the functions $\xi$, $\eta$, one decomposes the function $H$ into its linear factors
% [sic] The print has A_2 y in the second factor. Algebraically it must be A_1 y, as in the first factor, because A_1^2-4A_0A_2=-3D. The print is preserved.
\[
4\,A_0\,H = \bigl(2\,A_0\,x + A_1\,y + \sqrt{-3\,D}\,y\bigr)\,\bigl(2\,A_0\,x + A_2\,y - \sqrt{-3\,D}\,y\bigr).
\]
Then $\xi$, $\eta$ differ from these two factors of $H$ only by a constant factor each. We therefore set, if we denote two constant factors by $h$, $k$,
\weq{13}{
\begin{aligned}
2\,\xi &= h\,\bigl(2\,A_0\,x + A_1\,y - \sqrt{-3\,D}\,y\bigr) \\
2\,\eta &= k\,\bigl(2\,A_0\,x + A_1\,y + \sqrt{-3\,D}\,y\bigr),
\end{aligned}
}
from which, by multiplication, with regard to the first two equations (10), there results
\weq{14}{3\,h\,k\,A_0\,\sqrt[3]{-D} = 1,}
and the equations (11) give, by comparison of the coefficients of $x^3$,
\weq{15}{
\begin{aligned}
6\,h^3\,A_0^3\,\sqrt{-3\,D} &= +\,q_0 + 3\,\sqrt{-3\,D}\,a_0 \\
6\,k^3\,A_0^3\,\sqrt{-3\,D} &= -\,q_0 + 3\,\sqrt{-3\,D}\,a_0,
\end{aligned}
}
where $q_0$ is the coefficient of $x^3$ in $Q$, that is
\weq{16}{q_0 = 27\,a_0^2\,a_3 - 9\,a_0\,a_1\,a_2 + 2\,a_1^3}
.
